\textbf{Phần 1. Trắc nghiệm lựa chọn}
\setcounterpageref{ex}{0}

\begin{ex} %[5D1N4]
Phân số nào dưới đây là phân số thập phân?
\choice
{$\dfrac{9}{20}$}
{\True $\dfrac{63}{100}$}
{$\dfrac{100}{59}$}
{$\dfrac{15}{30}$}
\loigiai{
Phân số thập phân là phân số có mẫu số là $10; 100; 1\,000;\dots$. Trong các phân số đã cho, $\dfrac{63}{100}$ có mẫu số là $100$ nên là phân số thập phân.\\
\textbf{Đáp án của câu hỏi là B}
}
\end{ex}

\begin{ex} %[5D1N4]
Trong các phân số $\dfrac{7}{10}; \dfrac{9}{20}; \dfrac{178}{100}; \dfrac{16}{1\,000}; \dfrac{100}{59}$, có bao nhiêu phân số thập phân?
\choice
{$2$}
{$4$}
{$5$}
{\True $3$}
\loigiai{
Các phân số thập phân là: $\dfrac{7}{10}$ (mẫu số là 10), $\dfrac{178}{100}$ (mẫu số là 100), $\dfrac{16}{1\,000}$ (mẫu số là 1 000). Có tất cả $3$ phân số thập phân.\\
\textbf{Đáp án của câu hỏi là D}
}
\end{ex}

\begin{ex} %[5D1H4]
Viết phân số $\dfrac{3}{5}$ dưới dạng phân số thập phân có mẫu số là $10$, ta được
\choice
{$\dfrac{3}{10}$}
{$\dfrac{9}{10}$}
{\True $\dfrac{6}{10}$}
{$\dfrac{15}{10}$}
\loigiai{
Nhân cả tử số và mẫu số với $2$: $\dfrac{3}{5} = \dfrac{3 \times 2}{5 \times 2} = \dfrac{6}{10}$.\\
\textbf{Đáp án của câu hỏi là C}
}
\end{ex}

\begin{ex} %[5D1H4]
Chuyển phân số $\dfrac{11}{25}$ thành phân số thập phân có mẫu số là $100$, ta được
\choice
{$\dfrac{22}{100}$}
{\True $\dfrac{44}{100}$}
{$\dfrac{33}{100}$}
{$\dfrac{55}{100}$}
\loigiai{
Nhân cả tử số và mẫu số với $4$: $\dfrac{11}{25} = \dfrac{11 \times 4}{25 \times 4} = \dfrac{44}{100}$.\\
\textbf{Đáp án của câu hỏi là B}
}
\end{ex}

\begin{ex} %[5D1H4]
Rút gọn phân số $\dfrac{25}{500}$ về phân số thập phân có mẫu số là $100$, ta được
\choice
{\True $\dfrac{5}{100}$}
{$\dfrac{25}{100}$}
{$\dfrac{1}{100}$}
{$\dfrac{15}{100}$}
\loigiai{
Chia cả tử số và mẫu số cho $5$: $\dfrac{25}{500} = \dfrac{25 : 5}{500 : 5} = \dfrac{5}{100}$.\\
\textbf{Đáp án của câu hỏi là A}
}
\end{ex}

\begin{ex} %[5D1H4]
Chuyển phân số $\dfrac{82}{2\,000}$ thành phân số thập phân có mẫu số là $1\,000$, ta được
\choice
{$\dfrac{82}{1\,000}$}
{$\dfrac{20}{1\,000}$}
{\True $\dfrac{41}{1\,000}$}
{$\dfrac{164}{1\,000}$}
\loigiai{
Chia cả tử số và mẫu số cho $2$: $\dfrac{82}{2\,000} = \dfrac{82 : 2}{2\,000 : 2} = \dfrac{41}{1\,000}$.\\
\textbf{Đáp án của câu hỏi là C}
}
\end{ex}

\begin{ex} %[5D1H4]
Phân số $\dfrac{66}{60}$ viết dưới dạng phân số thập phân là
\choice
{$\dfrac{6}{10}$}
{$\dfrac{66}{10}$}
{$\dfrac{33}{10}$}
{\True $\dfrac{11}{10}$}
\loigiai{
Chia cả tử số và mẫu số cho $6$: $\dfrac{66}{60} = \dfrac{66 : 6}{60 : 6} = \dfrac{11}{10}$.\\
\textbf{Đáp án của câu hỏi là D}
}
\end{ex}

\begin{ex} %[5D1H4]
Phân số $\dfrac{27}{300}$ viết dưới dạng phân số thập phân là
\choice
{\True $\dfrac{9}{100}$}
{$\dfrac{3}{100}$}
{$\dfrac{27}{100}$}
{$\dfrac{18}{100}$}
\loigiai{
Chia cả tử số và mẫu số cho $3$: $\dfrac{27}{300} = \dfrac{27 : 3}{300 : 3} = \dfrac{9}{100}$.\\
\textbf{Đáp án của câu hỏi là A}
}
\end{ex}

\begin{ex} %[5D1H4]
Phân số $\dfrac{3}{4}$ viết dưới dạng phân số thập phân có mẫu số là $100$ là
\choice
{$\dfrac{25}{100}$}
{\True $\dfrac{75}{100}$}
{$\dfrac{50}{100}$}
{$\dfrac{30}{100}$}
\loigiai{
Nhân cả tử số và mẫu số với $25$: $\dfrac{3}{4} = \dfrac{3 \times 25}{4 \times 25} = \dfrac{75}{100}$.\\
\textbf{Đáp án của câu hỏi là B}
}
\end{ex}

\begin{ex} %[5D1H4]
Chuyển phân số $\dfrac{31}{125}$ thành phân số thập phân có mẫu số là $1\,000$, ta được
\choice
{$\dfrac{124}{1\,000}$}
{$\dfrac{62}{1\,000}$}
{$\dfrac{310}{1\,000}$}
{\True $\dfrac{248}{1\,000}$}
\loigiai{
Nhân cả tử số và mẫu số với $8$: $\dfrac{31}{125} = \dfrac{31 \times 8}{125 \times 8} = \dfrac{248}{1\,000}$.\\
\textbf{Đáp án của câu hỏi là D}
}
\end{ex}

\begin{ex} %[5D1H4]
Điểm chính giữa giữa hai phân số $\dfrac{3}{10}$ và $\dfrac{5}{10}$ trên tia số biểu diễn phân số thập phân nào?
\choice
{\True $\dfrac{4}{10}$}
{$\dfrac{2}{10}$}
{$\dfrac{6}{10}$}
{$\dfrac{8}{10}$}
\loigiai{
Trên tia số cách đều $\dfrac{1}{10}$, điểm nằm giữa $\dfrac{3}{10}$ và $\dfrac{5}{10}$ là $\dfrac{4}{10}$.\\
\textbf{Đáp án của câu hỏi là A}
}
\end{ex}

\begin{ex} %[5D1V4]
Phân số nào sau đây \textbf{không thể} chuyển thành phân số thập phân (có mẫu số là lũy thừa của $10$)?
\choice
{$\dfrac{7}{20}$}
{$\dfrac{9}{25}$}
{\True $\dfrac{5}{7}$}
{$\dfrac{13}{50}$}
\loigiai{
Mẫu số của các phân số thập phân chỉ có ước nguyên tố là $2$ và $5$. Phân số $\dfrac{5}{7}$ có mẫu số là $7$ nên không thể quy đồng về mẫu $10, 100, 1\,000,\dots$\\
\textbf{Đáp án của câu hỏi là C}
}
\end{ex}

\textbf{Phần 2. Trắc nghiệm đúng sai}
\setcounterpageref{ex}{0}

\begin{ex} %[5D1H4]
Cho các phân số: $A = \dfrac{11}{25}$, $B = \dfrac{25}{500}$, $C = \dfrac{66}{60}$, $D = \dfrac{31}{125}$. Xét tính đúng sai của các khẳng định sau:
\choiceTF
{\True Phân số $A$ được viết thành phân số thập phân là $\dfrac{44}{100}$.}
{Phân số $B$ viết thành phân số thập phân có mẫu số là $100$ là $\dfrac{25}{100}$.}
{\True Phân số $C$ viết thành phân số thập phân là $\dfrac{11}{10}$.}
{\True Phân số $D$ viết thành phân số thập phân có mẫu số $1\,000$ là $\dfrac{248}{1\,000}$.}
\loigiai{
\begin{itemize}
    \item $A = \dfrac{11 \times 4}{25 \times 4} = \dfrac{44}{100}$.
    \item $B = \dfrac{25 : 5}{500 : 5} = \dfrac{5}{100} \neq \dfrac{25}{100}$.
    \item $C = \dfrac{66 : 6}{60 : 6} = \dfrac{11}{10}$.
    \item $D = \dfrac{31 \times 8}{125 \times 8} = \dfrac{248}{1\,000}$.
\end{itemize}
\begin{itemize}
    \item a) \textbf{Đúng}.
    \item b) \textbf{Sai}: $B = \dfrac{5}{100}$.
    \item c) \textbf{Đúng}.
    \item d) \textbf{Đúng}.
\end{itemize}
}
\end{ex}

\begin{ex} %[5D1H4]
Xét tính đúng sai của các khẳng định sau về phân số thập phân:
\choiceTF
{\True Phân số $\dfrac{178}{100}$ là một phân số thập phân lớn hơn $1$.}
{Mọi phân số có tử số là $10, 100, 1\,000$ đều là phân số thập phân.}
{\True Phân số $\dfrac{3}{4}$ có thể chuyển thành phân số thập phân có mẫu số là $100$ với tử số bằng $75$.}
{Tổng của hai phân số thập phân $\dfrac{3}{10}$ và $\dfrac{57}{100}$ bằng $\dfrac{60}{100}$.}
\loigiai{
\begin{itemize}
    \item Phân số $\dfrac{178}{100}$ có mẫu số là $100$ và tử số $178 > 100$ nên là phân số thập phân lớn hơn $1$.
    \item Phân số thập phân yêu cầu \textbf{mẫu số} phải là $10, 100, 1\,000,\dots$, tử số là số tự nhiên bất kì. (Ví dụ $\dfrac{100}{59}$ không phải phân số thập phân).
    \item $\dfrac{3}{4} = \dfrac{3 \times 25}{4 \times 25} = \dfrac{75}{100}$.
    \item $\dfrac{3}{10} + \dfrac{57}{100} = \dfrac{30}{100} + \dfrac{57}{100} = \dfrac{87}{100} \neq \dfrac{60}{100}$.
\end{itemize}
\begin{itemize}
    \item a) \textbf{Đúng}.
    \item b) \textbf{Sai}: Phân số thập phân có mẫu số là $10, 100, 1\,000,\dots$
    \item c) \textbf{Đúng}.
    \item d) \textbf{Sai}: Tổng bằng $\dfrac{87}{100}$.
\end{itemize}
}
\end{ex}

\textbf{Phần 3. Trả lời ngắn}
\setcounterpageref{ex}{0}

\begin{ex} %[5D1H4]
Viết phân số $\dfrac{3}{5}$ dưới dạng phân số thập phân có mẫu số là $10$. Hỏi tử số của phân số thập phân đó là bao nhiêu?
\shortans{6}
\loigiai{
\textbf{Đáp số:} $6$\\
Ta có: $\dfrac{3}{5} = \dfrac{3 \times 2}{5 \times 2} = \dfrac{6}{10}$. Tử số là $6$.
}
\end{ex}

\begin{ex} %[5D1H4]
Chuyển phân số $\dfrac{11}{25}$ thành phân số thập phân có mẫu số là $100$. Tử số của phân số thập phân đó bằng bao nhiêu?
\shortans{44}
\loigiai{
\textbf{Đáp số:} $44$\\
Ta có: $\dfrac{11}{25} = \dfrac{11 \times 4}{25 \times 4} = \dfrac{44}{100}$. Tử số là $44$.
}
\end{ex}

\begin{ex} %[5D1H4]
Viết phân số $\dfrac{27}{300}$ dưới dạng phân số thập phân có mẫu số là $100$. Tìm tử số của phân số đó.
\shortans{9}
\loigiai{
\textbf{Đáp số:} $9$\\
Ta có: $\dfrac{27}{300} = \dfrac{27 : 3}{300 : 3} = \dfrac{9}{100}$. Tử số là $9$.
}
\end{ex}

\begin{ex} %[5D1H4]
Viết phân số $\dfrac{31}{125}$ dưới dạng phân số thập phân có mẫu số là $1\,000$. Tìm tử số của phân số nhận được.
\shortans{248}
\loigiai{
\textbf{Đáp số:} $248$\\
Ta có: $\dfrac{31}{125} = \dfrac{31 \times 8}{125 \times 8} = \dfrac{248}{1\,000}$. Tử số là $248$.
}
\end{ex}

\begin{ex} %[5D1H4]
Rút gọn phân số $\dfrac{82}{2\,000}$ thành phân số thập phân có mẫu số là $1\,000$. Hỏi tử số của phân số thu được là bao nhiêu?
\shortans{41}
\loigiai{
\textbf{Đáp số:} $41$\\
$\dfrac{82}{2\,000} = \dfrac{82 : 2}{2\,000 : 2} = \dfrac{41}{1\,000}$. Tử số là $41$.
}
\end{ex}

\begin{ex} %[5D1V4]
Tìm số tự nhiên $x$ sao cho $\dfrac{x}{50}$ được viết dưới dạng phân số thập phân có mẫu số $100$ là $\dfrac{76}{100}$.
\shortans{38}
\loigiai{
\textbf{Đáp số:} $38$\\
Ta có: $\dfrac{x}{50} = \dfrac{2 \times x}{100} = \dfrac{76}{100} \Rightarrow 2 \times x = 76 \Rightarrow x = 38$.
}
\end{ex}
